\chapter{CONCLUSIONES Y RECOMENDACIONES}
\label{cap:conclusiones}

\section{Conclusiones del diseño}

La máquina diseñada es técnicamente viable en el nivel de ingeniería preliminar. Sostiene, centra, traslada e inmoviliza un rodamiento de 3032 kg con las barras y las uniones del bastidor verificadas por debajo de su capacidad, y en el modelo de mojado aplica el fluido a sus cuatro superficies sin zonas secas. El diseño está verificado por simulación y cálculo, no validado con ensayos: la remoción de grasa y la seguridad en operación dependen de pruebas que quedan fuera del alcance de la tesis. Su inversión estimada, de clase 4, es de S/~319\,095 sin IGV; con doce rodamientos al año, se recupera en cinco años si cada rodamiento aporta un beneficio de S/~7015, mucho más que los S/~885 de mano de obra que reemplaza.

El primer objetivo se cumple. Se definieron quince especificaciones a partir de la geometría y la masa del rodamiento, las necesidades de los usuarios y las restricciones del taller, y doce restricciones con su origen documental.

El segundo objetivo se cumple en el nivel conceptual. La descomposición funcional y la matriz morfológica generaron cuatro arquitecturas que integran soporte, traslado, aplicación del fluido y recogida de residuos, y los filtros de restricciones las admitieron a todas, algunas de forma condicionada.

El tercer objetivo se cumple. La matriz de decisión seleccionó la arquitectura con el rodamiento y la mesa inmóviles durante el lavado y el aplicador en movimiento, con 3,63 puntos frente a 3,35 de la segunda opción. La misma lógica seleccionó el centrador de brazo pivotante y el traslado por doble cadena con volante, que obtuvo 4,25 puntos frente a 3,70 de la versión motorizada.

El cuarto objetivo se cumple. El bastidor de siete perfiles W6$\times$9 tiene una razón demanda/capacidad máxima de 0,82, gobernada por el desvío lateral de las ruedas, y sus diez uniones a momento quedan en 0,58 como máximo. La rueda de 200 mm cumple con 0,89 en servicio y 0,43 en estado último, por lo que no hace falta la de 250 mm. La red de 24 boquillas requiere 221,07 L/min con presiones de 4,00 a 4,13 bar, y la bomba Grundfos CRN 15-3 con variador cubre los 53,39 m del cálculo y los 51,0 a 51,3 m del trazado modelado. El distribuidor giratorio, su eje y su cubo quedan con una razón máxima de 0,93 en un caso accidental con el rotor detenido; en ese mismo caso, la unión del transversal con el cubo fluye localmente y queda por reforzar o por acotar el golpe de diseño.

El quinto objetivo se cumple: la evaluación frente a los requisitos obligó a rediseñar. El error de colocación, que en la segunda entrega llegaba a 6,10 mm, baja a 4,09 y 4,24 mm al limitar la fuerza del empujador a 1,95 kN y a 3,35 mm con alas continuas en los nudos. La cara inferior, que recibía cuatro veces menos dosis que las demás superficies, alcanza 19,99 kg\,m$^{-2}$ por vuelta en su zona exterior con el par de boquillas de abanico estrecho, y el tiempo de lavado estimado de esa cara pasa de 36,7 a 14,7 min; con ese criterio, el tiempo de la máquina, 17,6 min, lo fija ahora la cara superior. Ese cambio tiene un costo aceptado: el diámetro exterior pierde uniformidad en su franja baja, aunque sigue mojado por completo. El lavado lleva la dosis requerida a las superficies exteriores, que son las que se inspeccionan visualmente con el rodamiento armado, pero no alcanza su interior: solo el 24\,\% de rodillos y pistas tiene línea de vista con las boquillas, de modo que su limpieza corresponde a la etapa con el rodamiento desarmado. La masa real del rotor, 231 kg, reemplazó a la envolvente de 180 kg en la verificación del accionamiento.

De las quince especificaciones, catorce se satisfacen en los modelos y en los cálculos; la velocidad del distribuidor es un valor de diseño que solo puede medirse en el equipo. Ninguna de las doce restricciones se incumple, pero solo la técnica se satisface sin condiciones: las demás dependen de datos del taller, de la losa, de la implementación de los controles o de ensayos. El compromiso central del proyecto fue aceptar un distribuidor giratorio más complejo para mantener inmóviles la pieza y la mesa durante el lavado, lo que simplifica el soporte y la seguridad.

El uso de estándares fue determinante en varias decisiones. AISC 360-16 fijó la verificación de las barras y las uniones soldadas del bastidor, y la guía de torsión del AISC explicó por qué las uniones gobiernan el centrado. El ajuste H7/g6 de ISO 286-2 acotó los juegos del centrado, EN 1991-3 dio la fuerza lateral de guiado que hoy gobierna el bastidor y ASME B31.3 se usó para la tubería del distribuidor. ISO 12100 y la serie EN 12921 estructuraron el análisis de riesgos, y la evaluación del empuje según ISO 11228-2 y la cifra guía de la HSE llevó a descartar el empuje directo de la mesa. ACI 318-19 y la norma E.030 se emplearon en los anclajes y en la revisión sísmica. La acreditación de conformidad con estas normas corresponde a la etapa de fabricación.

\section{Recomendaciones}

\paragraph{Manufactura y escala.} Antes de fabricar, se recomienda completar los planos de pieza de las barras, de la cabina y del bloque portaboquillas, calificar los procedimientos de soldadura del taller y confirmar la losa real, cuyo espesor y resistencia se supusieron. La obra civil del pozo, del piso húmedo y de los pedestales debe diseñarse con esos datos. Para rodamientos de otras dimensiones, la arquitectura puede escalarse ajustando la posición de los apoyos sectoriales, la U y la lanza, pero cada tamaño exige repetir la verificación del bastidor y la simulación del mojado.

\paragraph{Materiales y equipos.} Se recomienda obtener del fabricante la curva certificada de la bomba en su punto de operación y su cotización con el motor de 4 kW, la capacidad de la rueda a baja velocidad, la compatibilidad del revestimiento de los apoyos con el detergente a 62~$^\circ$C y la resistencia de la cadena inoxidable. El calentamiento del baño debe diseñarse a partir del balance de la Sección~\ref{sec:calentamiento}: las resistencias, el aislamiento del tanque y el control de temperatura quedaron como estimación. La motorización del traslado puede montarse sobre la misma transmisión, cuyo reductor admite un motor con brida; la potencia de la Ecuación~\eqref{eq:dis-potencia-traslado} es el punto de partida, una vez conocido el rendimiento de la transmisión. Si se quiere extender el lavado a la hilera inferior de rodillos, conviene evaluar la boquilla adicional bajo el aro interior de la Sección~\ref{sec:anx-cavidad}, comprobando antes su espacio libre entre la lanza y las bases de los apoyos.

\paragraph{Seguridad.} Se recomienda diseñar el control eléctrico de la máquina con un controlador programable, con los enclavamientos que impiden la aspersión sin la puerta cerrada y la mesa bloqueada, y el traslado sin el distribuidor detenido a 0$^\circ$, y fijar su nivel de desempeño de seguridad. Deben calibrarse el disparo del limitador de fuerza del centrador y el par del limitador del distribuidor, y medirse la repetibilidad de las referencias y el juego del pasador de bloqueo. La valoración de riesgos debe revisarse con el responsable de seguridad del taller.

\paragraph{Sostenibilidad.} Se recomienda que una etapa posterior mida la grasa retirada por superficie, empezando por la cara inferior, para fijar el tiempo de ciclo que limpia sin exceso de fluido ni energía. El secado del rodamiento después del enjuague y el tratamiento del efluente quedan por diseñar: el efluente debe caracterizarse para definir su gestión, y el secado debe dimensionarse con la temperatura final del rodamiento y el agua retenida. Con cotizaciones formales de las diez partidas de mayor costo y con los datos de solvente, residuos y parada del molino que tiene el cliente, podrá confirmarse si la máquina se justifica económicamente.
